\begin{abcliste}
\abc A proton (charge $e$) gains $e U$ through the voltage $U$:
\begin{align}
 U &= \UF = \frac{\EkO}{e} = \result{\UP{6}{3}}
\end{align}

\abc Compare the kinetic energy with the rest energy of the proton,
\begin{align}
 E_0 &= \EzF = \Ez \approx \EzP{6}{4}
\end{align}
$E_\mathrm{kin}/E_0$ is far from small: relativity is needed. The total energy is $E = E_0 + E_\mathrm{kin} = \gamma E_0$:
\begin{align}
 \gamma &= \gmF = 1 + \frac{\EkO}{\EzP{6}{4}} = \gm
\end{align}
From $\gamma = 1/\sqrt{1 - v^2/c^2}$:
\begin{align}
 \frac{v}{c} &= \btF = \bt \approx \result{\btP{0}{2}}
\end{align}
The protons fly at about \btP{0}{2} times the speed of light. The classical formula $\frac{1}{2} m v^2 = E_\mathrm{kin}$ would give $v/c = \bcF = \bcP{0}{2}$, too much.
\end{abcliste}
